\subsection{The positive LM4 towers: scalar stability and the retained base}

This audit recovers the positive constructions underlying LM4 A.76,
A.79 and A.81 from the complete checkpoint, including the original
finite witnesses.  Its conditional inputs are stated explicitly.
In particular, a proof extending specified data over the congruence
base does not independently identify those data with the historical
arithmetic CM frames.  That last identification is independently
proved in E.95--E.97 below; the displayed extension arguments
retain their explicit inputs so that the dependence is visible.

Write \(\mathfrak B=C_2\amalg C_3\), with generators \(s,r\),
\(\tau=sr\), and
\[
 c(s)=s,\qquad c(r)=sr^{-1}s,\qquad c(\tau)=\tau^{-1}.
\]
All commutator subgroups and generated subgroups below are closed.
For a normal subgroup \(W\), put \(W_0=W\) and
\(W_{j+1}=[W_j,\mathfrak B]\).

\begin{result}{E.90}{Proved}{Exact scalar stabilization over a relatively perfect base}
Let \(R\triangleleft\mathfrak B\), \(R=[R,\mathfrak B]\), and
\(C=\mathfrak B/R\).  Suppose that \(W\triangleleft\mathfrak B\)
is contained in \(R\), that \(A=R/W\) is finite abelian, and that
\[
 \mathfrak B/W=A\rtimes C,\qquad
 H^1(C,A^\vee)=0,\qquad
 A^\vee=\operatorname{Hom}(A,\mathbb Q/\mathbb Z).
                                                        \tag{E.90.1}
\]
Let \(H\) be the finite action image on \(A\).  For every prime \(p\)
with \(A_p\ne0\), assume that \(H\) has a normal subgroup \(J_p\)
of order prime to \(p\) with \((A_p)_{J_p}=0\).  Set
\[
 M_p=(\bigwedge\nolimits^2_{\mathbb Z}A_p)_{O^p(H)},
 \quad M=\prod_pM_p,\quad
 m=\max\{1,e_2+1,2e_3+1\},
                                                        \tag{E.90.2}
\]
where \(p^{e_p}\) is the exponent of \(A_p\).
Then the actual source quotients satisfy
\[
 W/W_m\simeq M,\qquad
 (W_j/W_m)_p=I_p^jM_p,\qquad W_{m+1}=W_m.                \tag{E.90.3}
\]
Here \(I_p\) is the augmentation ideal of \(H/O^p(H)\).
If \(H\) acts trivially on \(M\), then already
\[
 W/[W,\mathfrak B]\simeq M,\qquad
 [[W,\mathfrak B],\mathfrak B]=[W,\mathfrak B].
                                                        \tag{E.90.4}
\]
The hypotheses hold automatically when
\(R=R_{\rm nil}(K)\) for a closed normal \(K\triangleleft\mathfrak B\),
provided \(A=R/W\) is finite abelian.  These conclusions identify
source quotients; they impose no norm equation on an automorphism.
\end{result}
\begin{proof}
The image of the pronilpotent residual in every finite quotient is
the stable lower-central subgroup.  Compactness of the fibres proves
this assertion and then \([R_{\rm nil}(K),K]=R_{\rm nil}(K)\).
For the automatic case put \(N=K/R\), a normal pronilpotent subgroup
of \(C\).  The coinvariants of \(A_p\) under \(N\) vanish.
After taking coinvariants under the prime-to-\(p\) factor \(N_{p'}\),
the remaining action is a finite \(p\)-group action.  A finite
abelian \(p\)-group with zero coinvariants under a \(p\)-group is
zero: its augmentation filtration terminates.  Consequently
\((A_p)_{N_{p'}}=0\).  Averaging and Hochschild--Serre give
\[
 H^i(C,A_p)=H^i(C,A_p^\vee)=0\quad(i\ge0).
                                                        \tag{E.90.5}
\]
In particular, the finite-kernel extension splits.  Such an
extension has a continuous set section: an open subgroup avoiding
the finite kernel gives local inverses on a finite clopen partition.
Thus the usual continuous extension class is defined.  The image
of \(N_{p'}\) supplies the required \(J_p\).

We prove the upper bound under (E.90.1).  At each step
\(\mathfrak B/W_i\to\mathfrak B/W_{i-1}\) has globally central
kernel.  Its \(R\)-image has abelianization \(A\).
Indeed, in a finite quotient retaining \(A\), a possible extra
abelianization kernel \(Z_0\) is a trivial \(C\)-module.  Dualize
\(0\to Z_0\to B\to A\to0\).
The vanishing of \(H^1(C,A^\vee)\) would lift any nonzero character
of \(Z_0\) to a \(C\)-invariant character of \(B\), whereas
\([R,\mathfrak B]=R\) excludes such a character.
Finite quotients retaining \(A\) are cofinal, so the assertion
holds before taking a finite quotient as well.

Inductively, the derived \(R\)-image is a finite quotient of
\(\bigwedge^2 A\), and the \(p\)-part of its action factors through
\(H/O^p(H)\).  At the next step a possible third commutator is an
equivariant pairing from the previous derived group tensored with
\(A\) into a globally trivial group.  The first factor is
\(O^p(H)\)-trivial and \((A_p)_{O^p(H)}=0\), so that pairing
vanishes.  The new group is again class two.  The automorphisms
fixing both successive derived subquotients form a \(p\)-group,
being a subgroup of a group \(\operatorname{Hom}(D,Z)\).
This proves the action assertion at every step.
The maximal pro-\(p\) quotient of \(\mathfrak B\) is \(C_2\),
\(C_3\), or \(1\), according as \(p=2\), \(3\), or another prime.
Thus
\[
 I_2^2=-2I_2,\qquad I_3^3=-3I_3(1+I_3),
\]
and \(I_2^{e_2+1}=I_3^{2e_3+1}=0\) on the indicated groups.
This gives stabilization and a natural surjection
\(M\to W/W_m\).

For equality, the actual source must attain that upper bound.
Put \(T_p=\operatorname{Hom}(A_p,M_p)\) and
\(X_p=\operatorname{Ext}^1_{\mathbb Z}(A_p,M_p)\).
Their \(J_p\)-invariants vanish.  For \(X_p\), an invariant
abelian extension admits a strict \(J_p\)-action by coprime
cohomology; its averaging idempotent retracts onto \(M_p\),
so that extension splits.  Hence all the \(H\)-cohomology of
\(T_p\) and \(X_p\) vanishes.
The commutator exact sequence
\[
 0\to X_p\to H^2(A_p,M_p)\to
       \operatorname{Hom}(\bigwedge^2A_p,M_p)\to0
                                                        \tag{E.90.6}
\]
therefore lifts the canonical alternating form to an invariant
extension class.  This gives a finite class-two group \(P_p\)
with \(P_p^{\rm ab}=A_p\), \(P_p'=M_p\), and a strict \(H\)-action.
The last assertion uses \(H^2(H,T_p)=0\).

In \((\prod_pP_p)\rtimes C\), lift the actual source elements
\(s,r\).  Correct their squares or cubes centrally whenever the
relevant integer is a unit.  The remaining \(2\)-primary
\(s\)-square is fixed by all of \(H\), because its action factors
through the cyclic quotient generated by \(s\).  Pulling back
\(C_{2^{j+1}}\to C_2\) along \(C\to C_6\to C_2\), and pushing
the central generator to the negative of that square, cancels it.
The analogous \(3\)-primary carry cancels the \(r\)-cube.
Here \(C\to C_6\) exists since \(R\subset\mathfrak B'\).
The resulting actual elements satisfy \(S^2=R^3=1\), hence
define a source homomorphism.  Its \(R\)-image maps onto \(A\);
commutators then attain \(M\), so the map is onto.
The images of \(W_j\) are \(I^jM\).  Comparison with the upper
bound proves (E.90.3), including its commutator identifications.
If \(H\) acts trivially, the image of \(W_1\) is already zero,
and the same upper-bound argument gives (E.90.4).
\end{proof}

\paragraph{The precise base input for the following two results.}
Let \(K_{\rm cong}\) be the congruence kernel of \(\mathfrak B\),
put \(R=R_{\rm nil}(K_{\rm cong})\), and \(C=\mathfrak B/R\).
Assume given a positive peripheral automorphism \(h_C\), with
negative involution \(b_C=h_Cc_C\), literal cusp inversion and
norm-one \(s/r\)-frames.  Require its action on the finite
head \(H\) of E.59 to be the identity, and its frame projections
there to be the stored pure projections \(\beta_0\) and \(s\).
These are the exact data used by the finite models.
Their historical production is the A.31/A.33 base assertion,
independently established with these exact head projections in
E.95--E.97.  It is not imported as an unproved local lemma.

\begin{result}{E.91}{Proved}{The supplied base data extend through the complete integral scalar tower}
Under the stated base input, the original A.56 witness has an
ordinary full positive realization, retaining its literal finite
frames.  The actual integral tower has nested relatively perfect
kernels \(U_n\), with residual abelianizations and centres
\[
 \begin{aligned}
 A_n&=\mathcal O_n^{36\cdot4}\oplus V^4,
       &\mathcal O_n&=\operatorname{GR}(2^n,3),\\
 M_n&\simeq\mathcal O_n^6\oplus\mathbb F_8^{30},
       &|M_n|&=2^{18n+90}.
 \end{aligned}                                          \tag{E.91.1}
\]
Its actions and two norm-one frames can be chosen compatibly,
with literal cusp, and realized by one ordinary full positive
automorphism.  The newer central frame coordinates may change.
The kernel \(U_\infty=\bigcap_nU_n\) is relatively perfect.
This proves the positive extension mechanism of A.56/A.62
from its specified base input.
\end{result}
\begin{proof}
We first check that the relevant action head is an actual quotient
of \(C\).  It is \(H=\mathbb F_7^{13}\rtimes Q_6\), where
\(Q_6=S_3\times A_4\) is congruence and its relation module is
the actual graph module of E.59.  Its congruence-kernel image
is all of \(\mathbb F_7^{13}\): every solvable quotient of a
congruence quotient has no \(C_7\) composition factor.
Indeed, the congruence groups
at primes at least five are perfect, and the solvable parts
at two and three have only those two primes.
Since this kernel image is abelian, \(R\) is killed.
The normal prime-to-two factor of \(K_{\rm cong}/R\)
maps onto the seven-primary module, giving the acyclicity
used in E.90.

The complete eight source cocycles of E.59 define an actual
affine quotient \(V^8\rtimes C\).
They remain independent on \(R\) by \(H^1(C,V)=0\);
simplicity with endomorphism field \(k=\mathbb F_8\) proves
surjectivity.  The opposite-weight calculation gives
\[
 (\bigwedge\nolimits^2_{\mathbb F_2}V^8)_{O^2(H)}
       =\operatorname{Sym}^2_k(k^8),
                                                        \tag{E.91.2}
\]
with trivial \(H\)-action.  The actual class-two law, including
its orientation and source carries, is the one verified in E.59.
E.90 identifies it with the first scalar refinement and proves
its kernel relatively perfect.  The recovered verifier
\path{code/lm4_source_verifiers/work_A/check_Q6_final_frames.py}
was executed on the original witness.  It verifies the action,
literal cusp, both source conjugacy equations and both norm-one
frames.  Thus E.57 gives the asserted ordinary full realization.
No finite search is used to prove absence of a lift.

For the integral tower use the unramified cubic ring
\(\mathcal O_\infty/\mathbb Z_2\) and the same monomial matrices,
with Teichm\"uller seventh roots.  The \(s\)-permutation has
four fixed lines and sixteen pairs.  Its invariant module has
rank twenty.  The \(r\)-norm kernel has rank twenty-four.
After normalizing \(d(s)=0\), quotienting by the
twenty-dimensional coboundary space leaves a free module of
rank four.  Cusp restriction is an isomorphism
\[
 \ker\bigl(H^1(\mathfrak B,\mathcal O_n^{36})
              \to H^1(\langle s\rangle,\mathcal O_n^{36})\bigr)
       \simeq \mathcal O_n^{36}/(\tau-1)\mathcal O_n^{36}
       \simeq\mathcal O_n^4.                             \tag{E.91.3}
\]
These statements follow integrally from unit minors modulo two.
The restored integral certificate verifies those minors,
the four admissible cusp cycles, and the actual coefficient
involution.  In particular the latter is the identity on cusp
invariants and coinvariants: its scalar on a fixed admissible
cycle is a seventh root reducing to one, hence is one.

Choose the four free cusp cocycles compatibly over
\(\mathcal O_\infty\), while retaining the four binary
\(s\)-classes separately.  Reduction gives the actual
surjection \(\mathfrak B\to A_n\rtimes C\), since its residual
image is onto modulo two and finite Nakayama applies.
Keeping the binary classes separately is necessary;
the entire torsion cohomology need not lift inside one
free coefficient module.
Write \(W_n\) for the source kernels and
\(U_n=[W_n,\mathfrak B]\).

Only opposite odd weights survive in the exterior square.
Their transporter exchanges the two weight lines, imposing
minus transpose on the integral multiplicity module
\(T_n=\mathcal O_n^4\oplus k^4\).  Thus the centre is
\[
 (T_n\otimes_{\mathcal O_n}T_n)/
       \langle u\otimes v+v\otimes u\rangle
       \simeq\mathcal O_n^6\oplus k^{30}.                \tag{E.91.4}
\]
There are six free off-diagonal coordinates, twenty-two other
off-diagonal coordinates, and eight order-two diagonal coordinates.
The three disjoint Frobenius weight families give precisely
the underlying binary group in (E.91.1), not three additional
copies of it.  The head action is trivial.  E.90 proves
\([U_n,\mathfrak B]=U_n\).  Reduction of the actual cocycles
gives \(W_{n+1}\subset W_n\), hence \(U_{n+1}\subset U_n\).

Choose a full lift of the original finite datum just constructed,
fixing \(\tau\), and denote its negative companion by \(b=hc\).
On the free cusp cohomology \(h\) acts as one and \(c\) as minus
one.  For an actual cocycle \(d\), unique vectors \(w,v\) satisfy
\[
 d(hg)=d(g)+(1-\rho(g))w,\qquad
 d(cg)=-C_0d(g)+(1-\rho(cg))v.                            \tag{E.91.5}
\]
Here \(C_0\) is the coefficient involution.
Literal cusp fixation gives \(w\in L_n^\tau\), so \(C_0w=w\);
applying \(c\) twice gives \(C_0v=v\).
Consequently \(d(b^2g)=d(g)\).
Together with the retained binary action, this proves preservation
of every \(W_n\) and square one on its affine quotient.
On the scalar refinement, \(b^2(g)g^{-1}\) is a central
homomorphism.  It kills \(\tau\), which generates
\(\mathfrak B^{\rm ab}=C_6\), so it vanishes.

The all-exponent frame step is more than a calculation modulo four.
No \(s\)-orbit exchanges the two weights of an opposite pair.
For the selected weight \(\lambda\), a contrary occurrence would give
\[
 q^{-1}sq\in\{\beta_0,\beta_0s\}\subset S_3\times A_4.
                                                        \tag{E.91.6}
\]
The first right-hand element has trivial \(S_3\)-component,
and the second has trivial \(A_4\)-component, whereas \(s\)
has neither.  Hence neither is conjugate to \(s\).
The recovered \path{check_A62_pair_orientation.py} independently
checks all thirty-six weights: there are two fixed unordered
pairs and eight two-element orbits of pairs.
Choose an \(s\)-invariant orientation.  The monomial \(s\)-action
then has no orientation carry at any exponent.
Every \(s\)-fixed \(a\in2L_n^4\), with central coordinate zero,
centralizes the actual source \(s\): its mixed commutators
with the binary coordinates vanish and the source has zero
free cusp coordinate.

Here are the literal affine frame equations in one free coefficient
factor.  Write the negative action as
\[
 b(p,g)=\bigl(-C_0p+(1-\rho(bg))q,bg\bigr),\qquad C_0q=q.
                                                        \tag{E.91.7}
\]
For the prescribed base \(s\)-frame \(\beta\), put
\(D_s=\rho(\beta)C_0\).  Substitution in the conjugacy and
normalizer-square equations for \((p,\beta)\) gives exactly
\[
 \rho(s)(p-q)=p-q,\qquad D_s(p-q)=p-q.                   \tag{E.91.8}
\]
The two operators commute and have square one.  For \(r\), first
replace the cocycle by \(d'=d-(1-\rho)a\) so that \(d'(r)=0\);
this is possible since three is a unit.  Then
\(q'=q-a-C_0a\) is \(C_0\)-fixed.  For the prescribed
negative \(r\)-conjugator \(J\), set
\(R'=\rho(b(r))\) and \(D_r=\rho(J)C_0\).  The equations are
\[
 R'(p-q')=p-q',\qquad D_r(p-q')=p-q',
 \quad (R')^3=D_r^2=1,\quad D_rR'D_r=(R')^{-1}.          \tag{E.91.9}
\]
Undoing the coboundary change restores the original source and
the same literal old \(J\).  Thus these are invariant modules
for monomial quotients of \(C_2^2\) and \(S_3\), respectively.
All stabilizer characters into seventh roots are trivial,
so an orbitwise rescaling makes permutation modules.
Their invariant modules reduce surjectively at every exponent.
The \(s\)-centralizer assertion just proved lifts the resulting
affine correction on the correct side of the full conjugator.
For \(r\), any central conjugacy discrepancy has both
two-power order and cube one, so is zero.
The two normalizer squares now lie in
\(\ker(M_n\to M_1)=2\mathcal O_n^6\).
The negative action is the identity on these free centre coordinates.
Independent central halves therefore kill both squares.
This changes only newer central frame coordinates.
The original-law A.61 verifier was rerun and checks its exact
modulo-four witness; the preceding argument proves every exponent.

For a finite set of levels, the largest level and E.57 give a
full positive realization with the required quotient action
and frames.  Closedness in the compact space of full conjugators,
including generation in every finite quotient, gives one
simultaneous realization.  No surjectivity between frame-solution
sets has been used.
Finally, for every finite quotient \(q\),
\[
 q(U_\infty)=\bigcap_nq(U_n);
\]
the descending finite images stabilize and equal their commutators
with \(q(\mathfrak B)\).  Finite separation proves
\([U_\infty,\mathfrak B]=U_\infty\).
\end{proof}

\begin{result}{E.92}{Proved}{The six-cup tower is relatively perfect and controls all central iterates}
Under the same base input, the changed stored A.64 datum extends
through the actual six-cup extension of every integral level,
with literal cusp, square-one negative action and both norm-one
frames.  Its finite-level and inverse-limit source kernels are
relatively perfect.  One ordinary full positive automorphism
realizes all these data.

Consequently the exact A.79 datum, from this base, extends
compatibly to every actual quotient
\(\mathfrak B/N_j\), \(N_{j+1}=[N_j,\mathfrak B]\).
The original A.76 datum has the asserted ordinary positive
realization and first-central compatibility by E.91.
These statements include the positive mechanisms in the historical
claims; their exact base marking is supplied by E.97.
The tower over \(C\) alone is not cofinal.
\end{result}
\begin{proof}
Use the four normalized binary cusp cocycles \(a_1,\ldots,a_4\),
and \(T=(V^{(2)})^6\).  The actual group law is
\[
 (a,z)(b,w)=
 \bigl(a+b,z+w+(a_i*b_j)_{i<j}\bigr),
 \qquad a,b\in V^4.                                    \tag{E.92.1}
\]
Multiplication is coordinatewise at each physical weight;
the head action on \(T\) has squared scalars.
The derived group and Frattini subgroup are both \(T\).
The actual source has zero \(s\)-cocycle, and the \(r\)-image
is obtained by conjugating the pure order-three element.
It follows from \(H^1(C,V)=0\) and the Frattini criterion
that the source surjects onto this group over \(C\).

Pairing it with the actual affine integral quotient gives the
complete \(V^8\) Frattini image.  It is therefore again onto.
Its image of \(W_\infty\) is \(T\); since \([T,H]=T\), its
image of \(U_\infty\supset[W_\infty,\mathfrak B]\) is all of \(T\).
Thus the six-cup extension of \(D_\infty=\mathfrak B/U_\infty\)
is the full actual fibre product, and the same holds at finite levels.
Its residual \(Q_n\) over \(C\) is class two, with
\[
 Q_n^{\rm ab}=A_n,\qquad Q_n'=M_n\oplus T.               \tag{E.92.2}
\]
The two summands are attained independently: they have trivial
and nontrivial odd weights, respectively.

The original A.64 tuple was recovered unchanged and checked by
\path{code/lm4_source_verifiers/work_A/check_A64_cross_cup_witness.py}.
The affine witness hash begins \texttt{4135fabf};
the six-cup witness hash begins \texttt{f5d10fa9}.
The manifest records both complete hashes.
The verifier substitutes the stored cocycle shears into the
original multiplication, verifies the source orders, canonical
\(c\), literal cusp, negative square, both frame projections,
both conjugacy equations and both normalizer norms.
It neither searches for nor modifies a witness.
The new affine choice has zero free cusp coordinates in the
difference between its \(s\)-frame and negative inner term.
These zeros lift integrally by zero.
The other free inner coordinates are cusp invariant and lift
in the integral invariant module of E.91.
The central-half corrections there have zero \(T\)-coordinate.
Consequently the exact six-cup action and frames extend through
all integral levels.

It remains to prove source scalar stability, rather than assume it.
Let \(W'_n\) be the source kernel of this extension.
In a finite quotient of \(\mathfrak B/[W'_n,\mathfrak B]\)
retaining \(Q_n\), let \(\widetilde Q\) be the \(R\)-image
and \(Z\) its globally central kernel over \(Q_n\).
Such finite quotients are cofinal.
The dual-character argument of E.90 uses
\(H^1(C,A_n^\vee)=0\) and yields
\[
 \widetilde Q^{\rm ab}=A_n,\qquad Z\subset\widetilde Q'.
                                                        \tag{E.92.3}
\]
Because \(Q_n\) is class two, \(\gamma_4(\widetilde Q)=1\).
The possible third commutator is an equivariant pairing
\[
 (M_n\oplus T)\otimes A_n
          \longrightarrow\gamma_3(\widetilde Q)\subset Z.
                                                        \tag{E.92.4}
\]
The \(M_n\)-part vanishes by averaging over the normal odd subgroup,
since \(M_n,Z\) are trivial there and \(A_n\) has no coinvariants.

The \(T\)-part factors through \(A_n/2A_n=V^8\).
After scalar extension to \(k\), the additive binary representation
has the three disjoint weight sets
\(\mathcal S,2\mathcal S,4\mathcal S\), each closed under negation.
The cup bracket is zero between different Frobenius forms.
Within one form, two vectors of the same weight \(\lambda\)
attain every one of the six \(T\)-coordinates of weight \(2\lambda\).
Only a vector of weight \(-2\lambda\) can pair with this
to give a globally trivial third commutator.
It belongs to a different Frobenius form.
For such \(a,b,c\), Jacobi gives
\[
 [[a,b],c]=-[[b,c],a]-[[c,a],b]=0.                        \tag{E.92.5}
\]
The two right-hand brackets have no \(T\)-component;
any old \(M_n\)-component has zero further commutator.
The argument is well defined modulo two: changing a representative
by twice an element only changes an old \(M_n\)-value.
Faithful scalar extension proves that all of (E.92.4) vanishes.

Now \(\widetilde Q\) is class two and its derived subgroup
is a quotient of \(\bigwedge^2_{\mathbb Z}A_n\).
In particular \(Z\) is a two-group.
Take invariants under the original normal odd subgroup \(J\).
Averaging gives
\[
 E_0=(\widetilde Q')^J\twoheadrightarrow M_n,\qquad
 \ker(E_0\to M_n)=Z.                                    \tag{E.92.6}
\]
The head fixes both subquotients, so its action on \(E_0\)
has two-group image and factors through \(H/O^2(H)\).
The exterior-square map therefore factors through its
coinvariants \(M_n\).
Thus the finite group \(E_0\) is a quotient of \(M_n\)
and surjects onto \(M_n\); equality of orders forces \(Z=0\).
This proves \([W'_n,\mathfrak B]=W'_n\).
The finite-image stabilization argument of E.91 gives the same
identity for \(W'_\infty=\bigcap_nW'_n\).
E.57 now supplies the full ordinary positive realization.

For completeness, any previously retained finite family can be
included under its actual hypotheses: its kernel must be relatively
perfect, its residual abelianization over \(C\) must have no
simple constituent of type \(V\) or \(V^\vee\), and its actions
and frames must agree over \(C\).
Equivariant Goursat then gives the full fibre product, because
a common nontrivial two-group quotient would have a common simple
abelian quotient.  Five-term inflation--restriction gives the same
cohomological vanishing over that retained base.
The preceding proofs apply unchanged and finite intersections
retain any such family simultaneously.
The assertion that every older A.47--A.52 object has those exact
specified data is not inferred merely from its historical label.

Let \(U\) be the relatively perfect source kernel just constructed
and \(N_0\) the kernel of its old finite quotient.  Since
\(U\subset N_0\), induction gives
\[
 U=[U,\mathfrak B]\subset[N_j,\mathfrak B]=N_{j+1}.
                                                        \tag{E.92.7}
\]
Therefore all the actual central iterates are quotients of this
one fixed source quotient and inherit its compatible action and
frames.  They are finite.  If \(N\) has index \(d\), transfer
shows that every invariant character of \(N\) is annihilated
by \(6d\), since \(\mathfrak B^{\rm ab}=C_6\).
The group \(N/[N,\mathfrak B]\) is finitely generated of
bounded exponent, hence finite; induction applies.

Finally, the tower over \(C\) has pro-two residual group.
Every nonabelian composition factor of a finite quotient of \(C\)
is \(\operatorname{PSL}_2(\mathbb F_q)\) for a prime \(q\);
its pronilpotent residual adds only abelian factors.
The new pro-two extension changes none of these nonabelian factors.
The \(A_7\)-quotient in E.56 consequently does not factor through
this tower, and its common source kernel maps onto \(A_7\).
This noncofinality assertion concerns the tower over \(C\);
it is not asserted for an arbitrary additional family of perfect
residual quotients.
\end{proof}

\begin{result}{E.93}{Proved}{The fresh-prime construction admits the required anti-symplectic coefficient action}
Let \(\mathfrak B\twoheadrightarrow Q\) be a finite quotient
in which \(s,r\) have exact orders two and three.  Put \(t=sr\).
Suppose \(b_0\) is an involution with \(b_0(t)=t^{-1}\) and
specified frames
\[
 b_0(s)=\beta s\beta^{-1},\quad b_0(r)=Jr^{-1}J^{-1},
 \quad b_0(\beta)=\beta^{-1},\quad b_0(J)=J^{-1}.
                                                        \tag{E.93.1}
\]
For infinitely many primes \(p\nmid2|Q|\), the actual relation
module \(M=H_1(\ker(\mathfrak B\to Q),\mathbb F_p)\) admits
a \(b_0\)-semilinear involution whose action on boundary homology
is the negative cusp permutation and whose action on compact
homology is anti-symplectic.
\end{result}
\begin{proof}
The kernel is the profinite completion of the torsion-free modular
surface subgroup, so its graph and boundary homology are the actual
ones.  Over \(\overline{\mathbb Q}\), let \(\mathscr M\) be
its homology, \(\mathscr P\) the boundary submodule, and
\(\mathscr H=\mathscr M/\mathscr P\) the compact homology.
The cellular and boundary sequences give
\[
 [\mathscr H]=2[1]+[kQ]-[k(Q/\langle s\rangle)]
             -[k(Q/\langle r\rangle)]-[k(Q/\langle t\rangle)].
                                                        \tag{E.93.2}
\]
The trivial multiplicity is zero; a nontrivial simple \(S\)
has multiplicity
\[
 m_S=\dim S-\dim S^s-\dim S^r-\dim S^t.                  \tag{E.93.3}
\]
Semisimplicity is valid in characteristic zero
\citep[Theorem~7.4, p.~103]{SymMilneGT2025}.
Duality and twisting by \(b_0\) preserve these multiplicities
and commute with each other.

On a four-element orbit
\(S,S^\vee,S^{b_0},(S^\vee)^{b_0}\), choose an intertwiner
on one pair, its negative adjoint on the dual pair, and inverses
on the return maps.  This is an anti-symplectic involution.
If \(S\ne S^\vee\) and twisting fixes both classes, rescale its
intertwiner \(A\) to \(A^2=1\), and use \(A\) on one block
and \(-(A^{-1})^\vee\) on the dual block.

If instead twisting exchanges \(S,S^\vee\), a nondegenerate
twisted form \(F\) satisfies
\[
 F(qx,b_0(q)y)=F(x,y).
                                                        \tag{E.93.4}
\]
Its transpose is \(F\) or \(-F\).  In the alternating case
the multiplicity \(m_S\) is even.
Indeed, the forms \(F(x,\beta y)\) and \(F(x,Jy)\) are
alternating by (E.93.1); \(s\) is an isometry for the former,
\(r\) is self-adjoint for the latter, and \(t\) is self-adjoint
for \(F\).  Their eigenvalue-one spaces are nondegenerate
alternating spaces.  Every dimension in (E.93.3) is even.
Tensoring \(F\) with an alternating multiplicity form in this
case, and a symmetric form otherwise, gives a symmetric twisted
form \(A\) on the full block \(X\).
The map
\[
 (x,\phi)\longmapsto(A^{-1}\phi,Ax)\quad\text{on }X\oplus X^\vee
                                                        \tag{E.93.5}
\]
is the required anti-symplectic involution.

For a self-dual \(S\) fixed by twisting, write the symplectic
block form as \(F\otimes G\), where \(F\) has symmetry
\(\varepsilon\) and \(G\) symmetry \(-\varepsilon\).
Rescale a twisted intertwiner \(A\) to \(A^2=1\).
Its multiplier on \(F\) is \(\nu=\pm1\).
If \(\nu=-1\), use \(A\otimes1\).
If \(\nu=1\), use \(A\otimes J_0\), where \(J_0\) is an
involutory anti-isometry of \(G\).
For alternating \(G\), a symplectic basis supplies \(J_0\).
For symmetric \(G\), its dimension is even: \(F\) is then
alternating and all the fixed-space dimensions in (E.93.3)
are even.  A hyperbolic basis supplies \(J_0\).
When twisting exchanges two self-dual classes, choose a twisted
anti-isometry and its inverse.  Over the algebraically closed
field, the required forms of the same symmetry and dimension
are isometric, with a scalar correction if necessary.
This exhausts the joint orbits.

Extend by the signed permutation on
\(\mathscr P=k[Q/\langle t\rangle]/k\), and choose a rational
\(Q\)-splitting of \(\mathscr M\to\mathscr H\).
All the resulting matrices are defined over a number field.
Pass to a finite Galois field containing their entries and avoid
their denominators and determinants.  Infinitely many rational
primes split completely in it
\citep[Theorem~8.31 and Corollary~8.32, p.~147]{MilneANT2020}.
Only existence is needed: for a monic primitive-element polynomial
\(f\), \(f(0)\ne0\), and a product \(m\) of finitely many bad primes,
\(f(f(0)mX)/f(0)\) is integral and congruent to one modulo \(m\).
A prime divisor of one of its values of absolute value greater
than one gives a new root of \(f\) modulo that prime.
Avoiding the discriminant gives an unramified degree-one prime,
and Galoisness makes it completely split.
Reducing the constructed matrices at such primes proves the result.
Boundary and compact integral homology are torsion-free, so the
excluded determinants ensure the actual source modules are retained.
\end{proof}

\begin{result}{E.94}{Proved}{First-central compatibility and ordinary realization at fresh primes}
Keep E.93's data, assume \(c\) preserves the kernel \(L\) of
\(\mathfrak B\to Q\), and suppose the same negative data extend
with both norm-one frames to the actual first central refinement
\[
 Q_1=\mathfrak B/[L,\mathfrak B].
 \]
For infinitely many primes \(p\nmid2|Q|\), put
\(N=\overline{L'L^p}\) and \(G_p=\mathfrak B/N\).
There is an involutory negative peripheral action on \(G_p\),
over \(b_0\), which extends with literal cusp and both frames
to \(\mathfrak B/[N,\mathfrak B]\).
Its positive companion has an ordinary full positive realization.

Applied to the A.79 data supplied from the base in E.92, this
proves the positive first-central and ordinary-realization clauses
of A.81.  E.59 supplies the separate obstruction to a full negative
involution, even with an unspecified cusp conjugator.
No statement about every central iterate at the new prime,
or about a particular numerically chosen prime, is made.
\end{result}
\begin{proof}
Let \(M=H_1(L,\mathbb F_p)\), let
\(P\subset M\) be its boundary module, and set \(H_c=M/P\).
The transfer argument in E.92 makes the central refinements
of both finite-index source kernels \(L,N\) finite.
The graph character
\([M]=[1]+[kQ]-[k(Q/\langle s\rangle)]-[k(Q/\langle r\rangle)]\)
has trivial multiplicity zero, so \(M_Q=0\).
Write \(W=(\bigwedge^2M)_Q\).
The \(p\)-primary part of the actual central refinement is
precisely
\[
 E_p=\mathcal E(M)\rtimes Q,\qquad
 (a,z)(b,w)=(a+b,z+w+\tfrac12(a\wedge b)_Q).
                                                        \tag{E.94.1}
\]
Here the source splitting \(G_p=M\rtimes Q\) is supplied
by coprime cohomology.

To prove the assertion about the actual quotient, let
\(\widetilde P\) be the preimage of \(M\) in its central
\(p\)-refinement.  A possible extra abelianization kernel
is \(Q\)-trivial.  Exactness of coinvariants, since
\(p\nmid|Q|\), would give a nonzero invariant character
of \(L\), contradicting \(M_Q=0\).
Thus \(\widetilde P^{\rm ab}=M\), and its central kernel
is its derived subgroup, a quotient of \(W\).
Its \(p\)-power map is additive for odd \(p\), equivariant
with trivial target, and hence zero.
The upper bound is attained by the actual source elements
\[
 S=(a_s,0,s),\qquad
 R=(a_r,-\tfrac16(a_r\wedge ra_r)_Q,r).
                                                        \tag{E.94.2}
\]
They satisfy \(S^2=R^3=1\).
The source \(L\)-image maps onto \(M\) and its commutators
attain \(W\); therefore it is the full group in (E.94.1).

An inner change of the whole marked model gives
\[
 \tau=(a_0,z,t),\qquad ta_0=a_0,\qquad
 \tau^m=(ma_0,mz,1),\quad m=\operatorname{ord}(t).
                                                        \tag{E.94.3}
\]
It transports both source elements together and does not
change their actual kernel.  The vector \(a_0=[\tau^m]/m\)
generates \(P\) as a \(Q\)-module.
Let \(B\) be the involution of E.93, so \(Ba_0=-a_0\),
and \(B_W\) its exterior-square action.
The homogeneous action on (E.94.1) is involutory and has
central cusp discrepancy \(z_0=(1+B_W)z\).

In \(W/(P\wedge M)_Q=(\bigwedge^2H_c)_Q\), the actual
surface relation gives
\[
 \Omega+|Q|\bar z=0,\qquad
 \Omega=\sum_i A_i\wedge B_i.                            \tag{E.94.4}
\]
There are \(|Q|/m\) positive boundaries, each with central
coordinate \(m\bar z\).  This explains the scalar and sign
in (E.94.4).
Anti-symplecticity sends \(\Omega\) to \(-\Omega\).
Since \(|Q|\) is invertible, \(z_0\) belongs to
\(I=(P\wedge M)_Q\).
The map
\[
 M^t\longrightarrow I,\qquad v\longmapsto(a_0\wedge v)_Q
                                                        \tag{E.94.5}
\]
is onto: move a translate of \(a_0\) back by coinvariants
and average the other vector over \(\langle t\rangle\).
It is equivariant for \(-B\) on its source and \(B_W\)
on its target.  Averaging over the group of order two
therefore gives \(v\in M^t\) with
\[
 Bv=-v,\qquad(a_0\wedge v)_Q=-z_0.
                                                        \tag{E.94.6}
\]
Compose the homogeneous action with conjugation by \((v,0,1)\).
Its square remains one and (E.94.6) makes its cusp literally
inverse.  Its action on \(G_p\) may change; that action was a
permitted choice.

The prescribed lower frames lift to \(G_p\).
One direct verification uses the extended module
\(M\rtimes(Q\rtimes\langle b_0\rangle)\):
the relevant order-four and order-six normalizer groups
are prime to \(p\), so averaging splits them and conjugates
their lifts to the actual order-two or order-three source
subgroups.  Equivalently, their frame discrepancies are
the usual coprime one-cocycles.
Any further lift of these frames to \(E_p\) has the exact
conjugacy equation: a central \(p\)-discrepancy with square
or cube one is zero.
The normalizer square is \(B_W\)-fixed in \(W\); its central
half removes it.  Both frames are now norm one.

Finally compare the old central primes.
Transfer proves that the kernel of \(Q_1\to Q\) has exponent
dividing \(6|Q|\), so \(|Q_1|\) is prime to \(p\).
For \(\ell\ne p\), inflation--restriction across \(N\lhd L\)
gives
\[
 \operatorname{Hom}_{\mathfrak B}(L,\mathbb Q_\ell/\mathbb Z_\ell)
 \simeq
 \operatorname{Hom}_{\mathfrak B}(N,\mathbb Q_\ell/\mathbb Z_\ell).
                                                        \tag{E.94.7}
\]
Positive cohomology of the finite \(p\)-group \(L/N\)
vanishes on these coefficients, and the isomorphism is
equivariant for \(\mathfrak B/L\).
Consequently the actual first refinement is the fibre product
\[
 \mathfrak B/[N,\mathfrak B]=E_p\times_QQ_1.              \tag{E.94.8}
\]
Its source image is the full product because the two residual
kernels have coprime orders.
The new action and frames pair with the assumed ones on \(Q_1\).
Their positive companion satisfies E.57, which supplies a full
ordinary positive lift inducing the specified \(G_p\)-action.
An inner correction in \(N\) fixes the full cusp literally.
The criterion retains the action on \(G_p\); it need not retain
the entire chosen action on its first refinement.
This is exactly the distinction required in A.81.
\end{proof}
